%% Chapter 9, Section 9.3 — Optimal Particle Filter
%% A First Course in Monte Carlo Methods
%% Sanz-Alonso & Al-Ghattas

\section{Optimal Particle Filter}
\label{sec:9.3}

The bootstrap particle filter propagates particles through the dynamics and then assigns
them weights according to the likelihood defined by the observation model. In contrast,
the optimal particle filter stems from a different conceptual decomposition of the
filtering step, where we first apply Bayes' formula and then propagate particles through a
Markov kernel which depends on the new observation.

\subsection{Derivation of the Optimal Filter Decomposition}
\label{sec:9.3.1}

Let $\mathbb{P}(\cdot|\cdot)$ denote conditional density. The following calculation
illustrates the decomposition of the filtering step that the optimal particle filter
approximates via particles:
\begin{align*}
  \mathbb{P}(X_{j+1} | Y_{1:j+1})
  &= \int_{\R^d} \mathbb{P}(X_j, X_{j+1} | Y_{1:j+1})\,dX_j \\
  &= \int_{\R^d} \mathbb{P}(X_{j+1} | X_j, Y_{1:j+1})\,\mathbb{P}(X_j | Y_{1:j+1})\,dX_j \\
  &= \int_{\R^d} \mathbb{P}(X_{j+1} | X_j, Y_{j+1})\,\mathbb{P}(X_j | Y_{1:j+1})\,dX_j \\
  &= \int_{\R^d} \mathbb{P}(X_{j+1} | X_j, Y_{j+1})\,
       \frac{\mathbb{P}(Y_{j+1} | X_j, Y_{1:j})}{\mathbb{P}(Y_{j+1} | Y_{1:j})}\,
       \mathbb{P}(X_j | Y_{1:j})\,dX_j \\
  &= \int_{\R^d} \mathbb{P}(X_{j+1} | X_j, Y_{j+1})\,
       \frac{\mathbb{P}(Y_{j+1} | X_j)}{\mathbb{P}(Y_{j+1} | Y_{1:j})}\,
       \mathbb{P}(X_j | Y_{1:j})\,dX_j \\
  &= \mathcal{P}_j^{\mathrm{opt}} \mathcal{A}_j^{\mathrm{opt}} \mathbb{P}(X_j | Y_{1:j}),
\end{align*}
where
\begin{align*}
  \mathcal{P}_j^{\mathrm{opt}} \pi(x_{j+1})
    &= \int_{\R^d} \mathbb{P}(x_{j+1} | X_j, Y_{j+1})\,\pi(X_j)\,dX_j, \\
  \mathcal{A}_j^{\mathrm{opt}} \pi(x_j)
    &\propto \mathbb{P}(Y_{j+1} | x_j)\,\pi(x_j).
\end{align*}

We have therefore derived the following decomposition of the filtering step:
\begin{equation}
  \label{eq:9.7}
  f_{j+1} = \mathcal{P}_j^{\mathrm{opt}} \mathcal{A}_j^{\mathrm{opt}} f_j.
\end{equation}

\subsection{Implementation}
\label{sec:9.3.2}

In general, it is not possible to implement the optimal particle filter in the fully
nonlinear setting due to two computational bottlenecks:

\begin{itemize}
  \item There may not be a closed formula for evaluating the likelihood
    $\mathbb{P}(Y_{j+1} | X_j)$, making unfeasible the computation of the particle
    weights.

  \item It may not be possible to sample from the Markov kernel
    $\mathbb{P}(X_{j+1} | X_j, Y_{j+1})$, making unfeasible the propagation of
    particles.
\end{itemize}

However, when the observation function $b(\cdot)$ is linear, i.e.\ $b(\cdot) = H\cdot$
for some matrix $H$, both bottlenecks may be overcome. We thus consider the following
setting, which arises in many applications:
\begin{alignat*}{3}
  X_{j+1} &= a(X_j) + \xi_{j+1}, &\qquad \xi_{j+1} &\sim \Normal(0, \Sigma), \\
  Y_{j+1} &= HX_{j+1} + \eta_{j+1}, &\qquad \eta_{j+1} &\sim \Normal(0, \Gamma),
\end{alignat*}
where $X_0$, $\{\xi_j\}$, and $\{\eta_j\}$ are all assumed to be independent. First,
note that
\[
  Y_{j+1} = Ha(X_j) + H\xi_{j+1} + \eta_{j+1},
\]
which implies that $\mathbb{P}(Y_{j+1} | X_j) = \Normal\!\left(Ha(X_j),\,
H\Sigma H^T + \Gamma\right)$. Second, note that
\[
  \mathbb{P}(X_{j+1} | X_j, Y_{j+1})
  \propto \exp\!\left(
    -\frac{1}{2}\bigl|Y_{j+1} - HX_{j+1}\bigr|_\Gamma^2
    - \frac{1}{2}\bigl|X_{j+1} - a(X_j)\bigr|_\Sigma^2
  \right),
\]
which implies via completion of the square that $\mathbb{P}(X_{j+1} | X_j, Y_{j+1}) =
\Normal(m_{j+1}, C)$, with
\begin{align*}
  m_{j+1} &= (I - KH)a(X_j) + KY_{j+1}, \\
  C &= (I - KH)\Sigma, \\
  K &= \Sigma H^T S^{-1}, \\
  S &= H\Sigma H^T + \Gamma.
\end{align*}
These formulas are similar to those arising in the Kalman filter algorithm, which you will
derive in Exercise~9.2. We are now ready to present a practical implementation of the
optimal particle filter with linear observations.

\begin{algorithm}[H]
  \caption{Optimal Particle Filter}
  \label{alg:9.2}
  \begin{algorithmic}[1]
    \Require Initial distribution $f_0$, observations $Y_{1:J}$, sample size $N$.
    \State \textbf{Initial sampling:} Sample $X_0^{(1)}, \ldots, X_0^{(N)} \iid f_0$
      so that $f_0^N = S^N f_0$.
    \State \textbf{Subsequent sampling:}
    \State For $j = 0, \ldots, J-1$, do for $1 \leq n \leq N$:
    \State Set
      \[
        \widehat{X}_{j+1}^{(n)} = (I - KH)a\!\left(X_j^{(n)}\right) + KY_{j+1}
          + \zeta_{j+1}^{(n)},
        \qquad
        \zeta_{j+1}^{(n)} \iid \Normal(0, C).
      \]
    \State Set
      \[
        \tilde{w}_{j+1}^{(n)} = \exp\!\left(
          -\frac{1}{2}\Bigl|Y_{j+1} - Ha\!\left(X_j^{(n)}\right)\Bigr|_S^2
        \right).
      \]
    \State Normalize the weights:
      \[
        w_{j+1}^{(n)} = \frac{\tilde{w}_{j+1}^{(n)}}{\displaystyle\sum_{n=1}^{N}
          \tilde{w}_{j+1}^{(n)}}.
      \]
    \State Draw $X_{j+1}^{(n)} \iid \sum_{n=1}^{N} w_{j+1}^{(n)}
      \delta\!\left(x - \widehat{X}_{j+1}^{(n)}\right)$.
    \State Set
      \[
        f_{j+1}^N(x) = \frac{1}{N} \sum_{n=1}^{N} \delta\!\left(x - X_{j+1}^{(n)}\right).
      \]
    \Ensure Approximation $f_J^N$ to the filtering distribution $f_J$.
  \end{algorithmic}
\end{algorithm}

\medskip
\noindent\fbox{\parbox{\dimexpr\textwidth-2\fboxsep-2\fboxrule\relax}{%
\textbf{Pros and Cons}: Similar to the bootstrap particle filter, the optimal particle
filter is convergent in the large $N$ limit, but also suffers from small effective sample
size in high-dimensional settings. The variance of the weights is smaller, however, for
the optimal particle filter than for the bootstrap particle filter. It is important to
notice that, in contrast to the bootstrap filter, the optimal filter relies on Gaussian
assumptions on the noise in the dynamics and observation models, and on linearity of the
observations. Thus, it can only be deployed on a smaller class of problems.
}}
